\section{Questão 12}
Uma fonte de informação gera símbolos aleatoriamente a partir de um alfabeto de quatro letras,
$\{a, b, c, d\}$,
com probabilidades 
$P(a) = 1/2$, 
$P(b) = 1/4$, 
$P(c) = 1/8$, 
$P(d) = 1/8$.
%
Suponha que as gerações de símbolos sejam independentes. 
Um esquema de codificação escreve esses símbolos em códigos binários da seguinte maneira:
%
\begin{table}[htbp!]
    \centering
    \begin{tabular}{c| cccc}
    \toprule
        Símbolo & $a$ & $b$ & $c$ & $d$ \\
    \midrule
        Código & 0 & 10 & 110 & 111 \\
    \bottomrule
    \end{tabular}
\end{table}

Seja $X$ a variável randômica denotando o comprimento do código em bits, ou seja, o número de bits da palavra digital.

\begin{multicols}{2}
\begin{itemize}
    \item [a)] Qual a faixa dinâmica de $X$?
    \item [b)] Encontre a probabilidade $P(X = 1)$,
    \item [c)] Determine $P(X = 2)$,
    \item [d)] Determine $P(X = 3)$.
    \item [e)] Encontre $P(X > 3)$.
    \item [f)] Especifique o tipo de $X$ e esboce a CDF.
    \item [g)] A partir da CDF, determine $P(X \leq 1)$,
    \item [h)] Análogo para $P(l < X \leq 2)$,
    \item [i)] Análogo para $P(l \leq X \leq 2)$, 
\end{itemize}
\end{multicols}

\respostas

\begin{itemize}
    \item [a)]
    Pela tabela, $X\in \{1, 2, 3\}$.
    
    \item [b)] 
    Como apenas a letra $a$ tem código com $X=1$ e $P(a)=1/2$, 
    segue que $P(X=1) = 1/2$,
    
    \item [c)]
    Como apenas a letra $b$ tem código com $X=2$ e $P(b)=1/4$, 
    segue que $P(X=2) = 1/4$,
    
    \item [d)] 
    Tanto a letra $c$ quanto $d$ satisfazem $X=3$. 
    Visto que a geração de símbolos é independente, podemos somar as probabilidades
    das letras serem geradas para obter $P(X=3) = P(c) + P(d) = 1/4$.
    
    \item [e)] 
    $X>3$ está fora da faixa dinâmica da variável randômica e portanto $P(X>3)=0$.
    
    \item [f)]
    Como $X$ é uma variável randômica discreta, a CDF é a soma cumulativa da PMF.
    Abaixo, são esboçados os gráficos das distribuições e a definição análitica de $F_X(x)$.
    
    \begin{minipage}{.32\linewidth}
        \centering 
        \begin{tikzpicture}[> = latex, scale=1.0]
            \draw [->] (-0.5, 0) -- (3.5, 0) node [right] {$x$};
            \draw [->] (0, 0) -- (0, 1.5) node [above] {$f_X(x)$};
            \draw 
                (1, 0) node [below] {$1$} -- (1, 0.5)
                (2, 0) node [below] {$2$} -- (2, 0.25)
                (3, 0) node [below] {$3$} -- (3, 0.25)
            ;
            \fill [azulunb]
                (0, 0) node [below, black] {$0$} circle (0.07)
                (1, 0.5) circle (0.07) node [above] {$1/2$}
                (2, 0.25) circle (0.07) node [above] {$1/4$}
                (3, 0.25) circle (0.07) node [above] {$1/4$}
            ;
        \end{tikzpicture}
    \end{minipage}
    %
    \begin{minipage}{.32\linewidth}
        \centering 
        \begin{tikzpicture}[> = latex, scale=1.0]
            \draw [->] (-0.5, 0) -- (3.5, 0) node [right] {$x$};
            \draw [->] (0, 0) -- (0, 1.5) node [above] {$F_X(x)$};
            \draw 
                (1, 0) node [below] {$1$} -- (1, 0.5)
                (2, 0) node [below] {$2$} -- (2, 0.75)
                (3, 0) node [below] {$3$} -- (3, 1)
            ;
            \fill [azulunb]
                (0, 0) node [below, black] {$0$} circle (0.07)
                (1, 0.5) circle (0.07) node [above] {$1/2$}
                (2, 0.75) circle (0.07) node [above] {$3/4$}
                (3, 1) circle (0.07) node [above] {$1$}
            ;
        \end{tikzpicture}
    \end{minipage}
    %
    \begin{minipage}{.32\linewidth}
        \begin{align}
            F_X(x) = (1/2)\ &\delta[t-1] + \nonumber\\
                     (3/4)\ &\delta[t-2] + \nonumber\\ 
                     &\delta[t-3] 
        \end{align}
    \end{minipage}
    
    \item [g), h), i)] Tem-se
    \begin{align}
        P(X \leq 1) &= F_X(1) = 1/2 \\
        P(1 < X \leq 2) &= F_X(2) - F_X(1) = 1/4 \\
        P(1 \leq X \leq 2) &= F_X(2) - F_X(0) = 3/4
    \end{align}
\end{itemize}